\section{Funding, Process, and Policy}
\label{sec:policy}

\subsection{Messaging, or, Who Pays For All This?}

\lesson{Exam administration has many potential stakeholders. Each success attracts support; cast a wide
net to bootstrap flexible funding, and message your successes.} 

We began with an initial internal ``microgrant'' of US\$10,000 from the
Division of Undergraduate Education, based on presentations we had
made to multi-department audiences (e.g.\ Deans and Chairs) about the
benefits of mastery learning and the key role of computer-based
testing in supporting it.  This grant supported some initial work
to convert lower-stakes content to PrairieLearn, such as homeworks
and remote quizzes,
to demonstrate what would be
possible in a CBTF setting.
We then made every effort to spread the word about the benefits of mastery
learning, understanding that \emph{different supporters need to hear
different messages.\/}  Faculty are concerned about pedagogy but also
about the workload and quality-of-life of their TAs.  The SAC is
concerned about the overwhelming number of students needing proctoring
that accounts for exam accommodations.  Deans,
Chairs, and Vice-whatevers are concerned about how best to deploy
limited academic staff budgets to high-leverage activities, and most
agree that proctoring isn't one.
As PIs on the research grant, we got ourselves invited everywhere we could using every connection we
had: faculty lunches in our
own and colleagues' departments, college-level meetings of Deans and Chairs,
external industrial advisory board meetings, teaching and learning seminars.
Each time we gave a presentation, we tried to show incremental but
positive results that would resonate with that stakeholder.
Evangelizing everywhere, we believed, would help ensure that when
an opportunity appeared,  potential benefactors would be positively
predisposed to support the project.

While we have no counterfactual, we believe the approach worked:
after a few months we attracted additional
seed funding of US\$75,000 from the College of Engineering, in
recognition of the fact that the ``first wave'' of courses to
potentially benefit from the CBTF would be homed there.
Later, on the success of that funding, we won two expansive grants
totaling about US\$1.5M (shared with two 
other institutions) covering research and content development related
to mastery learning as well as CBTF operations.

Critically, all of these pools of money had considerable flexibility
in how they were spent, and there was some ``informal'' accounting
within the institution as well.  For example, we did not have to pay
to use the shared lab infrastructure since it was already a ``common
good'' service, and their staff were willing to put in some  modest
extra effort to help us run them as part-time CBTF spaces
without line-item financial compensation.
Similarly, we were able to borrow some internal IT staff time for various
aspects of on-site hosting and proctor hiring/supervision,
as well as having explicit support in the research grants for IT staff
time to manage and pay for
CBTF cloud resources.
As a result of these visible successes, when a large dedicated space miraculously became available as part of
a building remodel, central campus committed over US\$1.5M to convert
and operate it as a permanent CBTF starting in Fall 2026.

\subsection{Institutional Exam Policies}

\lesson{Consider asking forgiveness, rather than permission, for
accidentally transgressing the letter of the law in good faith.}

Many institutions have regulations related to when and/or where students can
take exams, especially final exams.
Our courses are classified as either having a
``written final exam'' or not, and that designation requires committee
approval to modify.
A course with a written final is assigned a specific, non-negotiable time and room
during finals week, student conflicts be damned.
Must a CBTF final overlap that time slot?
What if  \emph{all\/}  your students demand to take your final at that
time, and there aren't enough CBTF seats?
Does a computer-administered exam even count as a ``written final''?  How many
angels can dance on the head of a pin?

The CBTF paradigm is so thoroughly foreign
compared to ``traditional'' exam administration that you may have no
choice but to break a few eggs when faced with such conundra.
At our institution, strong student and faculty voices, positive or
negative, carry significant weight; hence our focus on
spreading the word, collecting positive testimonials, and
gathering and acting on negative feedback.
To date, students both with and without accommodations in two dozen
courses overwhelmingly
appreciate the scheduling flexibility, a majority  prefer CBTF exams
to paper exams, and no instructor who has tried CBTF exams
wants to go back. 

\label{sec:cheating}
A stakeholder we didn't expect to engage was our summer sessions program.
One student in the pilot performed much worse on the in-CBTF final than their  remote
quiz scores would predict;
investigation revealed that they had been systematically cheating on
the remote quizzes by having someone else take them and then
synthetically combining screencasts and live video to fool the remote
proctors.
This tactic was thwarted on the final by the CBTF's much higher exam
security.
That humbling and discouraging experience led us to 
adopt a CBTF-only policy for \emph{all\/} quizzes in the course---a change we were 
happy to make as instructors,
but which upset our summer sessions program since a course that cannot be
taken fully online has smaller revenue potential.

